\section{A wild class on a mixed-characteristic open surface}
We prove Theorem~\ref{thm:wild} by reducing the question to the
first two graded pieces of a ramification filtration.

Write $A=W(k)[[t]]$, where $k=\overline{\F}_2$, and put
\[
 R_0=\Spec A[1/2],\quad V=\Spec A[1/(2t)],\quad
 B=k[[t]],\quad E=k((t)),\quad C=\Frac W(k).
\]
Let $D=(A_{(2)})^h$ and $L=\Frac D$. Then $D$ is an excellent
henselian DVR with residue field $E$. The coefficient sheaf
$\Lambda=\F_2$ is constant on all these schemes; on the
characteristic-zero schemes we identify it with every tensor
power of $\mu_2$.

\begin{proof}[Proof of Theorem~\ref{thm:wild}]
The regular parameters of $A$ are $2,t$. Complete noetherian
local rings are excellent. Its residue field is algebraically
closed, so the proposed residual bound equals two.
The cohomological-dimension theorems of Gabber--Orgogozo give
$\cd_2(R_0),\cd_2(V)\le2$, and $\cd_2(L)=2$
\cite[Theorems 1.2 and 6.1]{GabberOrgogozo}. More explicitly,
the equal-characteristic dimension invariant of $E=k((t))$
is one; adjoining the mixed-characteristic henselian valuation
raises it by one. Henselian invariance and the characteristic-two
field bound give $\cd_2(D)=\cd_2(E)\le1$.

The complement $z=U\setminus V$ is the point with residue field
$E$. Henselian excision and localization for $\Spec D$ give
$H^3_z(U,\Lambda)=H^2(L,\Lambda)$. Localization on $U$ thus gives
\begin{equation}\label{eq:wild-coker}
 H^2(V,\Lambda)\longrightarrow H^2(L,\Lambda)
     \longrightarrow H^3(U,\Lambda)\longrightarrow0.
\end{equation}
The same sequences with any two-primary torsion sheaf show
vanishing above degree three: the restrictions to $D,L,V$
vanish above degrees $1,2,2$, respectively. Hence $\cd_2(U)\le3$.

The point $R_0\setminus V$ is $(t)$, with residue field $C$.
Here two is invertible, so codimension-one purity applies.
Since $H^3(R_0,\Lambda)=0$, the residue sequence is exact:
\[
 H^2(R_0,\Lambda)\longrightarrow H^2(V,\Lambda)
     \longrightarrow H^1(C,\Lambda)\longrightarrow0.
\]
The class of $u\in C^\times/C^{\times2}$ is the residue of
$\{t,u\}$. Thus every class in the first image of
\eqref{eq:wild-coker} is of the form
\begin{equation}\label{eq:wild-image}
 a+\{t,u\},\qquad a\in H^2(R_0,\Lambda),\quad u\in C^\times.
\end{equation}

For $m\ge1$, let $U^mH^2(L,\Lambda)$ be generated by symbols
whose first entry belongs to $1+2^mD$, and let $U^0=H^2$.
The absolute ramification index is one and the critical index
is two. Kato's symbol calculation, in
\cite[Proposition 6.1.2(1),(3), p.~51]{Suzuki}, gives
\[
 \operatorname{gr}^0H^2(L,\Lambda)=E^\times/E^{\times2},
 \qquad \operatorname{gr}^1H^2(L,\Lambda)=\Omega_E^1.
\]
The maps are $v\mapsto\{\widetilde v,2\}$ in degree zero
and $x\,d\log y\mapsto\{1+2\widetilde x,\widetilde y\}$
in degree one. Suzuki's
Propositions 8.1.4 (pp.~64--65), 8.2.1(1),(3) (p.~65),
and 8.2.3 (p.~66) show that
$H^2(R_0,\Lambda)\to H^2(L,\Lambda)$ is injective and strict
for this filtration, with images
\[
 B^\times/B^{\times2}\subset E^\times/E^{\times2},
 \qquad \Omega_B^1\subset\Omega_E^1
\]
on these two graded pieces. The hypotheses hold because $A$
is regular, excellent and henselian, its residue field is perfect,
$-1\in A$ is the required second root of unity, and its unique
prime dividing two has regular quotient $B$.

Suppose a class \eqref{eq:wild-image} lies in $U^1$. Write
$u=2^eu_0$, with $u_0\in W(k)^\times$. Since $k$ is perfect,
divide by a square of a Teichm\"uller unit to make $u_0\in1+2W(k)$.
The degree-zero contribution of the symbol is $t^e$, whereas
that of $a$ is the class of a unit of $B$. Equality modulo
squares forces $e$ even, by the $t$-valuation. Remove that
square; the symbol is now in $U^1$, and strictness puts $a$
in $U^1$ as well. For $u=1+2b$, its degree-one contribution
is $\bar b\,d\log t$, with $\bar b\in k$. The whole degree-one
image is therefore contained in
\[
 \Omega_B^1+k\,d\log t=B\,dt+k\,t^{-1}dt.
\]
On the other hand, the class $c=\{1+2/t,t\}$ has image
$t^{-2}dt$. This does not belong to the displayed subspace:
the coefficient of $t^{-2}$ separates it. Since $k$ is perfect,
$t$ is a two-basis of both $B$ and $E$, so their absolute
differential modules are exactly $B\,dt$ and $E\,dt$; no
completed-differential convention is being used. Thus $c$ is
outside the image of $H^2(V,\Lambda)$. Its boundary in
\eqref{eq:wild-coker} is nonzero, proving the theorem.
\end{proof}
